% =====================================================================
%  Fräsen mit Seitenanschlag, Führungsschiene und Stangenzirkel.
%  Draufsicht, schematisch. Selbst gezeichnet.
% =====================================================================
\begin{tikzpicture}[x=1mm, y=1mm, font=\scriptsize]

  % Frästisch von oben mit Fräser: \fraestisch{(x,y)}{Radius}
  \newcommand{\fraestisch}[2]{%
    \draw[metall, fill=black!25, fill opacity=0.9] #1 circle (#2);
    \fill[white] #1 circle (3.6);
    \fill[black!65] #1 circle (2);}

  % ---------------- 1 Seitenanschlag -----------------------------------
  \feld{0}{0}{54}{66}{1\quad Seitenanschlag}
  \draw[werkstueck] (3,22) rectangle (51,54);
  \fill[nut] (3,36) rectangle (27,40);
  \draw[holzrand, dashed, line width=0.4pt] (27,36) -- (51,36) (27,40) -- (51,40);
  \draw[metall, fill=black!10] (20,19) rectangle (21.6,38) (32.4,19) rectangle (34,38);
  \draw[griff, rounded corners=0.8] (7,16) rectangle (47,22);
  \fraestisch{(27,38)}{9}
  \draw[hilfslinie] (37,38) -- (48,38);
  \draw[mass] (46,22) -- (46,38);
  \node[anchor=west] at (46.5,30) {\textbf{a}};
  \draw[vorschub] (14,10.5) -- (34,10.5);
  \node[anchor=west] at (35,10.5) {Vorschub};
  \node at (27,4) {Anschlag liegt an der Kante an};

  % ---------------- 2 Führungsschiene ----------------------------------
  \feld{58}{0}{54}{66}{2\quad Führungsschiene}
  \draw[werkstueck] (61,22) rectangle (109,54);
  \fill[nut] (61,28) rectangle (85,32);
  \draw[holzrand, dashed, line width=0.4pt] (85,28) -- (109,28) (85,32) -- (109,32);
  \draw[metall, fill=black!30] (59,42) rectangle (111,49);
  \draw[black!60, line width=0.6pt] (59,45.5) -- (111,45.5);
  \foreach \x in {64,102}
    \draw[griff, rounded corners=0.6] (\x,40) rectangle (\x+4,51);
  \draw[metall, fill=black!10] (80.2,34) rectangle (81.8,45) (88.2,34) rectangle (89.8,45);
  \draw[griff, rounded corners=0.6] (78,41) rectangle (92,50);
  \fraestisch{(85,30)}{9}
  \draw[hilfslinie] (88,32) -- (99,32);
  \draw[mass] (97.5,32) -- (97.5,42);
  \node[anchor=west, fill=holz, inner sep=0.5pt] at (98.2,37) {$\ge$\,5\,mm};
  \draw[vorschub] (70,10.5) -- (90,10.5);
  \node[anchor=west] at (91,10.5) {Vorschub};
  \node at (85,4) {Schiene mit Zwingen festspannen};

  % ---------------- 3 Stangenzirkel ------------------------------------
  \feld{116}{0}{54}{66}{3\quad Stangenzirkel}
  \draw[werkstueck] (119,10) rectangle (167,54);
  \coordinate (M) at (140,30);
  \draw[holzrand, dashed, line width=0.4pt] (M) circle (14) (M) circle (18);
  \fill[nut] (M) ++(-150:18) arc[start angle=-150, end angle=20, radius=18]
    -- ($(M)+(20:14)$) arc[start angle=20, end angle=-150, radius=14] -- cycle;
  \draw[metall, fill=black!10] (M) -- ++(110:1) -- ($(M)+(20:16)+(110:1)$)
    -- ($(M)+(20:16)+(-70:1)$) -- ($(M)+(-70:1)$) -- cycle;
  \fraestisch{($(M)+(20:16)$)}{7.5}
  \draw[mass] (M) -- ++(125:16);
  \node at ($(M)+(125:9)+(35:2.5)$) {\textbf{r}};
  \draw[black!85, fill=white, line width=0.6pt] (M) circle (1.4);
  \draw[black!85, line width=0.5pt] (M) ++(-1,0) -- ++(2,0) (M) ++(0,-1) -- ++(0,2);
  \node at (143,4) {Zirkelspitze im Mittelpunkt};
\end{tikzpicture}
